\section{Deep ensembles}
\label{sec:deep-ensembles}

The three preceding sections approximated a posterior over the weights in three
different ways: a Gaussian variational family, a Bernoulli one, and a local
Gaussian read off the curvature at a single maximum. Deep ensembles dispense
with the posterior altogether. A fixed number $M$ of networks of the same
architecture is trained independently, and their predictive distributions are
combined. Lakshminarayanan et al.\ proposed the method explicitly as an
alternative to Bayesian neural networks, on the grounds that the latter demand
substantial changes to the training pipeline and are slower to train than
standard networks~\cite{lakshminarayanan2017}. Nothing in the construction
refers to a prior, a posterior or a divergence. The method therefore enters this
comparison from the opposite direction to
Sections~\ref{sec:bayes-by-backprop}--\ref{sec:laplace-approximation}, and it
closes the chapter from the non-Bayesian side.

The motivation the authors give comes directly from
Section~\ref{sec:mc-dropout}. Dropout admits a second reading, as a combination
of an ensemble of networks with shared parameters, alongside the variational one
established by Gal and Ghahramani. Lakshminarayanan et al.\ argue that the
ensemble reading is the more plausible of the two whenever the dropout rate is
not tuned on the training data, because any sensible approximation to a
posterior must depend on that data~\cite{lakshminarayanan2017}. This is the
objection already recorded in Section~\ref{subsec:mc-dropout-limitations},
where the rate $p_i$ was seen to be chosen rather than inferred. Taken
seriously, it suggests examining ensembles on their own terms rather than as a
by-product of a variational argument.

% ---------------------------------------------------------------------
\subsection{Construction and training}
\label{subsec:ensembles-training}

The recipe consists of three elements: a proper scoring rule as the training
criterion, an ensemble, and adversarial training as an optional
addition~\cite{lakshminarayanan2017}.

A scoring rule assigns a numerical value to a predictive distribution and
rewards better calibrated predictions over worse ones. It is proper when the
true distribution attains the optimum, and only then. The log-likelihood is a
proper scoring rule, whereas the squared error of a point prediction says
nothing about predictive uncertainty. For regression, each member therefore
carries the two output heads of Section~\ref{subsec:aleatoric} and is trained by
minimising the Gaussian negative log-likelihood of
Equation~\eqref{eq:gaussian-nll}, under the heteroscedastic model of
Equation~\eqref{eq:heteroscedastic}. The variance head is passed through a
softplus transform to enforce positivity, and a small floor is added for
numerical stability~\cite{lakshminarayanan2017}.

This is a design decision rather than a formality, and the authors measured its
effect. A single network that learns its variance achieved a better test
log-likelihood than an ensemble of five, or of ten, networks trained on the
squared error and equipped with the empirical variance of their predictions. The
calibration figures make the gap concrete: on the \emph{Year Prediction MSD}
dataset, the nominal $80\%$ interval derived from that empirical variance
contained roughly $20\%$ of the test observations~\cite{lakshminarayanan2017}.
The common heuristic of ensembling point predictors and reading the spread as an
uncertainty estimate is thus not a cheap version of the method described here.
It is a different and markedly worse one.

Diversity between members is the second element, and it is obtained at no cost.
Random initialisation of the parameters, together with random shuffling of the
training points, was found sufficient in practice. Every member is trained on
the whole training set rather than on a bootstrap sample, because bagging
degraded performance in the reported experiments; a bootstrap sample contains on
average only $63.2\%$ unique points, and a base learner with many local optima
does not need the extra randomisation~\cite{lakshminarayanan2017}. The same
protocol is adopted in this work, which also keeps the training data identical
to that used by the other four methods.

Why so little randomisation suffices was examined later by Fort et
al.~\cite{fort2019}. Independently initialised networks reach entirely different
modes of the loss landscape. Functions collected along a single optimisation
trajectory, or sampled from a subspace around it, cluster within one mode in the
space of predictions, even when they differ considerably in weight space.
Low-loss paths connecting distinct modes do exist, but the functions along them
are not connected in prediction space. Measured on the diversity--accuracy
plane, the decorrelating power of random initialisation is not matched by
subspace sampling methods~\cite{fort2019}.

This is the mechanism behind the epistemic estimates of the method, and it
locates the contrast with the two preceding sections precisely. Bayes by
Backprop is unimodal because the divergence of
Equation~\eqref{eq:vi-objective} is mode-seeking and the family of
Equation~\eqref{eq:mean-field} ignores correlations, as set out in
Section~\ref{subsec:bbb-limitations}. The Laplace approximation is unimodal
because the expansion is taken at one maximum, as set out in
Section~\ref{subsec:laplace-limitations}. Fort et al.\ propose exactly this
difference as the explanation for the gap between theory and practice: scalable
variational methods concentrate on a single mode, while deep ensembles cover
several, and Bayesian neural networks underperform ensembles in practice
particularly under dataset shift~\cite{fort2019}.

The third element is adversarial training. Given an input $x$ with target $y$
and a loss $\ell(\theta, x, y)$, the fast gradient sign method produces a
perturbed example
%
\begin{equation}
    x' = x + \epsilon \operatorname{sign}
      \bigl( \nabla_{x}\, \ell(\theta, x, y) \bigr) ,
    \label{eq:fgsm}
\end{equation}
%
which is added to the training set with the original target. The perturbation
follows the direction in which the loss rises fastest, so the procedure is read
as a cheap way of smoothing the predictive distribution over a neighbourhood of
each training point, rather than along a randomly chosen
direction~\cite{lakshminarayanan2017}. The step is marked optional in the
original algorithm, and the reported evidence is mixed. On MNIST it improved
every metric at every ensemble size, on SVHN the benefit faded as $M$ grew, and
on the regression benchmarks the difference between an ensemble with and without
it lay within the error bars~\cite{lakshminarayanan2017}.

Adversarial training was excluded from this work, for two reasons. The reported
gain on regression is not distinguishable from noise, and the perturbation size
$\epsilon$ is an additional hyperparameter that would have to be set per
dataset. Tuning it for one method alone would make the comparison depend on the
effort spent on that method, which is the argument already used for the kernel
choice in Section~\ref{subsec:gp-kernel}. A second consideration is that
Equation~\eqref{eq:fgsm} alters the training set, so the ensemble would no
longer be trained on the same data as the other four methods.

% ---------------------------------------------------------------------
\subsection{Prediction and uncertainty decomposition}
\label{subsec:ensembles-prediction}

The trained members are treated as a uniformly weighted mixture,
%
\begin{equation}
    p\bigl(y^\ast \mid x^\ast\bigr)
    = \frac{1}{M} \sum_{m=1}^{M}
      p_{\theta_m}\bigl(y^\ast \mid x^\ast\bigr) ,
    \label{eq:ensemble-mixture}
\end{equation}
%
where $\theta_m$ denotes the parameters of the $m$-th network. For Gaussian
members this is a mixture of Gaussians. To keep quantiles and predictive
probabilities easy to compute, the mixture is approximated by a single Gaussian
carrying its mean and its variance~\cite{lakshminarayanan2017}.

Both moments are already available. The mean is the average of the member means,
that is Equation~\eqref{eq:bbb-predictive-mean} with the index over stochastic
passes replaced by an index over members. The variance is
Equation~\eqref{eq:bbb-variance-estimator} under the same substitution, and it
is identical in form to Equation~\eqref{eq:mc-dropout-heteroscedastic}. The
estimator is therefore the third appearance of one expression, and it is not
restated here. What differs between the three methods is the origin of the
samples, and nothing else. In Section~\ref{sec:bayes-by-backprop} they are drawn
from a learnt variational posterior, in Section~\ref{sec:mc-dropout} from a set
of Bernoulli masks, and here each sample is the endpoint of an independent
optimisation run. The samples of a deep ensemble are not drawn from any
distribution at all.

% Uwaga dla autora: lakshminarayanan2017 zapisuje wariancję mieszaniny
% w postaci (1/M) sum (sigma_m^2 + mu_m^2) - mu_*^2, czyli z dzielnikiem M,
% tak samo jak kendall2017 w Sekcji 3.3. Powyżej przyjęto dzielnik T-1
% z Równania (bbb-variance-estimator). To ta sama decyzja, o której mowa
% w notatce w rozdzial33.tex — zadeklarować ją raz w Rozdziale 4
% i stosować dla wszystkich metod.

The reading of the two terms carries over unchanged from
Equation~\eqref{eq:total-variance-2}. The first averages the noise predicted by
the individual networks and is the aleatoric component. The second measures how
far the members disagree about the mean and is the epistemic component. It
vanishes for $M = 1$, which restates the observation of
Section~\ref{subsec:mle-map} that a single point estimate reports no epistemic
uncertainty by assumption. Where the training data are dense, independently
trained networks agree and the term is small. Outside the training range the
members are constrained by nothing in common, they diverge, and the term grows.
The toy experiment of the original work shows precisely this: ensembling
improved the predictive distribution most as the test points moved away from the
observed data~\cite{lakshminarayanan2017}.

% [Miejsce na Rysunek 3.10: posterior deep ensemble na zbiorze load_sin.
%  Trening na [0, 6], ewaluacja na [-2, 10]. TE SAME osie i ta sama skala co
%  Rysunek 3.2 (GP), 3.5 (BBB), 3.6 (MC dropout) i 3.8 (Laplace) — komplet
%  pięciu rysunków na wspólnej skali domyka główny argument wizualny
%  rozdziału.
%  Sieć: M niezależnych MLP z src/models.py (dropout = 0), różne ziarna
%  inicjalizacji i różna kolejność danych, trening Gaussian NLL.
%  Predykcyjna średnia z Równania (bbb-predictive-mean) po składowych,
%  pasmo +-2 sigma z Równania (bbb-variance-estimator), punkty treningowe,
%  prawdziwa funkcja linią przerywaną. OBOWIĄZKOWO: M pojedynczych średnich
%  cienką linią — to jedyne miejsce w pracy, gdzie widać wprost rozbieżność
%  składowych poza zakresem treningowym, czyli mechanizm z fort2019.
%  Cel rysunku: pokazać, że pasmo rozchodzi się poza [0, 6] szerzej niż
%  u BBB i MC dropout, i porównać tempo tego wzrostu z GP.]

\begin{figure}[htbp]
    \centering
    \includegraphics[width=0.85\textwidth]{rodzial3_rys/img3_10.png}
    \caption{Predictive distribution of a deep ensemble on the synthetic sine
    dataset. The models are trained on the interval $[0, 6]$ and evaluated on
    $[-2, 10]$, with the axes matching those of
    Figures~\ref{fig:gp-posterior-sine}, \ref{fig:bbb-posterior-sine},
    \ref{fig:mc-dropout-posterior-sine} and~\ref{fig:laplace-posterior-sine} to
    allow a direct comparison. Thin lines show the predictive means of the
    individual members, which coincide within the training range and separate
    outside it. The band shows $\pm 2\sigma$ obtained from the mixture of
    Equation~\eqref{eq:ensemble-mixture}.}
    \label{fig:ensemble-posterior-sine}
\end{figure}

% [Miejsce na Rysunek 3.11: wpływ liczby składowych M na oszacowanie
%  niepewności. Zbiór load_sin, ta sama architektura co na Rysunku 3.10.
%  Średnia szerokość pasma +-2 sigma oraz RMSE w funkcji M = 1, 2, 3, 5, 10,
%  15, 20, każdy punkt uśredniony po kilku losowaniach zestawu ziaren,
%  z pasmem odchylenia standardowego między losowaniami. Rysunek jest
%  odpowiednikiem Rysunku 3.7 dla MC dropout i ma tę samą rolę: uzasadnić
%  liczbowo wartość M przyjętą w Rozdziale 4, zamiast przyjmować ją
%  arbitralnie. Punkt odniesienia z literatury: M = 5 jako wartość domyślna
%  u lakshminarayanan2017 i malejące przyrosty powyżej pięciu składowych
%  u ovadia2019.]

\begin{figure}[htbp]
    \centering
    \includegraphics[width=0.85\textwidth]{rodzial3_rys/img3_11.png}
    \caption{Quality of the deep ensemble uncertainty estimate as a function of
    the number of members $M$. The mean width of the $\pm 2\sigma$ band and the
    root mean squared error are shown for increasing $M$, each averaged over
    several independent sets of initialisation seeds. The gains diminish well
    before the largest value of $M$, which fixes the ensemble size used in the
    experiments of Chapter~\ref{ch:setup}.}
    \label{fig:ensemble-size}
\end{figure}

% [Miejsce na Rysunek 3.12 (opcjonalny): wariancja uczona kontra wariancja
%  empiryczna. Dwa panele na wspólnych osiach, oba na zbiorze load_sin:
%  po lewej M sieci trenowanych na MSE, z pasmem +-2 sigma liczonym wyłącznie
%  z empirycznego rozrzutu predykcji; po prawej ten sam ensemble trenowany
%  na Gaussian NLL, z pasmem z Równania (bbb-variance-estimator).
%  Cel rysunku: odtworzyć na własnych danych wynik lakshminarayanan2017
%  uzasadniający dwugłowicową architekturę — heurystyka „uśrednij predykcje
%  punktowe i weź rozrzut" wyraźnie zaniża niepewność. Jeśli argument
%  liczbowy w tekście wystarczy, rysunek można pominąć.]

% ---------------------------------------------------------------------
\paragraph{Computational cost and limitations.}
The cost of the method is the highest in this comparison, and it is also the
easiest to describe. Training is repeated $M$ times and prediction requires $M$
forward passes, so both scale linearly with the ensemble size. The saving
argument advanced in Section~\ref{subsec:mc-dropout-limitations} applies here
without the qualification attached to it there: the members are trained
independently and share nothing, so the whole procedure parallelises trivially
and the wall-clock cost need not grow at all~\cite{lakshminarayanan2017}. Memory
is the component that does not parallelise away. The ensemble holds $M$ complete
copies of the network, which is the only entry in the cost table of Ovadia et
al.\ whose storage grows with the number of replications; Monte Carlo dropout
adds nothing and mean-field variational inference doubles the
count~\cite{ovadia2019}. Where memory is the binding constraint, the ensemble
can be distilled into a single smaller model~\cite{lakshminarayanan2017}.

Two limitations follow from the construction. The diversity on which the
epistemic term depends is an empirical property of random initialisation, not a
guaranteed one; Fort et al.\ measured it rather than derived
it~\cite{fort2019}. Nothing in the procedure encodes beliefs held before the
data are seen, so the method offers no counterpart to the kernel of
Section~\ref{subsec:gp-kernel} or the prior precision of
Section~\ref{sec:laplace-approximation} through which such beliefs could be
stated or tuned. A third qualification is specific to the present work. The
strongest evidence for deep ensembles under distribution shift comes from
classification benchmarks~\cite{ovadia2019, wilson2020}, whereas the comparison
in Chapter~\ref{ch:results} is a regression one, so those conclusions are
carried over as an expectation rather than as an established result.

That evidence is nevertheless consistent. Ovadia et al.\ benchmarked the
methods of this chapter across image, text and categorical data under
increasing corruption of the inputs. Their summary is that the quality of
uncertainty degrades with shift for every method, that calibration on
independent and identically distributed test data does not carry over to
shifted data, and that deep ensembles perform best across most metrics while
proving the most robust to shift~\cite{ovadia2019}. The advantage was shown not
to be a matter of capacity, since widening the single-network baselines did not
reproduce it, and a modest ensemble of five members already captured most of the
available gain~\cite{ovadia2019}.

A closing observation returns the method to the framework of
Chapter~\ref{ch:methods}. Wilson and Izmailov argue that the distinguishing
property of a Bayesian approach is marginalisation rather than optimisation, and
that deep ensembles are not a competitor to Bayesian inference but a mechanism
for performing it~\cite{wilson2020}. Because a network is typically
underspecified by the data, many different parameter settings explain it
equally well, which is exactly the situation in which averaging over them
matters. By representing several basins of attraction, an ensemble can
approximate the integral of Equation~\eqref{eq:predictive-2} more faithfully
than a method that represents one basin accurately. Their experiment supports
the claim directly. Against a reference predictive distribution built from long
Hamiltonian Monte Carlo chains, the deep ensemble was closer to the reference
than factorised variational inference, which was strongly overconfident between
clusters of data points; the distance to the reference fell with the number of
ensemble members, while for the variational method it barely moved with the
number of samples~\cite{wilson2020}.

This is the note on which the chapter closes.
Section~\ref{subsec:predictive-distribution} stated that the five methods differ
only in how they approximate Equation~\eqref{eq:predictive-2}. The last of them
was introduced as an alternative to Bayesian inference and can be read as a
particularly effective instance of it, which suggests that the label attached to
a method orders it less usefully than the quality of the approximation it
delivers. Chapter~\ref{ch:results} is arranged to test that ordering
empirically, and the expectation attached to deep ensembles differs from those
recorded for the earlier methods. Prediction intervals should be wider than
those of Bayes by Backprop and Monte Carlo dropout, empirical coverage should
sit closer to the nominal level than for either, and the growth of the epistemic
term under extrapolation should be the strongest among the four neural methods,
with the Gaussian process of Section~\ref{sec:gaussian-processes} as the
reference. The PICP and MPIW metrics of Chapter~\ref{ch:setup}, together with
Figure~\ref{fig:ensemble-posterior-sine}, are reported to settle it.
